% संपूर्ण मराठी ग्रंथातील प्रकरण 18: प्रथम-क्रम तर्कशास्त्रापलीकडे
% हा संपादनयोग्य स्रोतखंड आहे. संकलनासाठी संपूर्ण स्रोतसंग्रहातील
% अक्षररूपे, पूर्वभाग, चिन्हव्याख्या आणि आकृती-स्रोत आवश्यक आहेत.
% मूळ: Open Logic Project; मजकूर आणि रूपांतर CC BY 4.0.
% यंत्रानुवाद, दुरुस्ती आणि तपासणी: OpenAI Codex —
% GPT-5.6 Sol आणि GPT-6 Sol, दोन्ही Ultra effort.
% दुरुस्ती-पुनरावलोकन: GPT-6.1 Sol, Ultra effort.
\chapter{प्रथम-क्रम तर्कशास्त्रापलीकडे}\label{fol:byd::chap}
\begin{editorial}
  जेरेमी अविगॅड यांच्या तर्कशास्त्रावरील टिपणांवरून रूपांतरित केलेला हा
  अध्याय इतर कोणकोणत्या तार्किक प्रणाली अस्तित्वात आहेत यांची अत्यंत संक्षिप्त
  झलक देतो. ज्या अभ्यासक्रमात या प्रणाली शिकवल्या जात नाहीत, त्यात पुढील
  अध्ययनासाठी विषय सुचवणारा अध्याय म्हणून तो अभिप्रेत आहे. येथे उल्लेखिलेल्या
  प्रत्येक विषयाला पुढे---अशी आशा आहे---ओपन लॉजिक प्रोजेक्टमध्ये स्वतंत्र
  भागाच्या पातळीवरील विवेचन मिळेल.
\end{editorial}









\section{आढावा}\label{fol:byd:int:sec}

प्रथम-क्रम तर्कशास्त्र हीच काही महत्त्वाची एकमेव तार्किक प्रणाली नाही:
प्रथम-क्रम तर्कशास्त्राचे अनेक विस्तार आणि अनेक भेद आहेत. सामान्यतः एखाद्या
तर्कशास्त्रात भाषेचे आकारिक विनिर्देशन, बहुतेकदा---पण नेहमीच नाही---एक
निगामी प्रणाली आणि बहुतेकदा---पण नेहमीच नाही---एक अभिप्रेत चिन्हार्थमीमांसा
यांचा समावेश असतो. पण या संज्ञेच्या तांत्रिक वापरामुळे एक उघड प्रश्न निर्माण
होतो: सहजप्रज्ञ किंवा तात्त्विक अर्थाने वापरल्या जाणाऱ्या ``तर्कशास्त्र'' या
शब्दाशी प्रथम-क्रम तर्कशास्त्र नसलेल्या तर्कशास्त्रांचा संबंध काय? खाली
वर्णिलेल्या सर्व प्रणाली कोणत्या ना कोणत्या प्रकारच्या युक्तिविचाराचे प्रतिमान
देण्यासाठी रचल्या आहेत; पण त्या तार्किक कशामुळे ठरतात, हे आपण सांगू शकतो का?

याची सोपी उत्तरे मिळत नाहीत. ``तर्कशास्त्र'' हा शब्द वेगवेगळ्या अर्थांनी आणि
वेगवेगळ्या संदर्भांत वापरला जातो; आणि ``सत्य'' या संकल्पनेप्रमाणेच या
संकल्पनेचेही अनेक तात्त्विक भूमिकांतून विश्लेषण झाले आहे. उदाहरणार्थ, कोणती
विधाने अनिवार्यतः सत्य, अनुभवपूर्व सत्य, अतार्किक पदांच्या अर्थनिर्धारणापासून
स्वतंत्रपणे सत्य, त्यांच्या आकारामुळे सत्य किंवा भाषिक संकेतामुळे सत्य आहेत
हे ठरवणे हे तार्किक युक्तिविचाराचे ध्येय आहे, असे कोणी मानू शकेल; आणि यांपैकी
प्रत्येक संकल्पना बरीच स्पष्ट करावी लागते. केवळ गणितात वापरल्या जाणाऱ्या
तर्कशास्त्रापुरते लक्ष मर्यादित केले, तरी त्याची व्याप्ती किती याविषयी फारसे
एकमत नाही. उदाहरणार्थ, फ्रेगे यांनी सुरू केलेल्या {\em तर्कमीमांसावादी}
परंपरेत रसेल आणि व्हाइटहेड यांनी \textit{प्रिन्सिपिया मॅथेमॅटिका} मध्ये
तर्कशास्त्राच्या पायावर गणित विकसित करण्याचा प्रयत्न केला. त्यांची तार्किक
प्रणाली खाली वर्णिलेल्या प्रणालीसारख्या उच्च-प्रकार तर्कशास्त्राचे एक रूप
होती. शेवटी त्यांना बहुतेक निकषांनुसार पूर्णतः तार्किक न वाटणारी स्वयंसिद्धके
समाविष्ट करावी लागली (विशेषतः अनंतत्वाचे स्वयंसिद्धक आणि न्यूनीकरणक्षमतेचे
स्वयंसिद्धक); तरीही उच्च-क्रम युक्तिविचाराची काही रूपे तार्किक म्हणून
स्वीकारली पाहिजेत, असे कोणी मानू शकेल. याउलट, ज्यांच्या सत्ताशास्त्रात
``विधानांना'' संवादाच्या वैध वस्तू म्हणून स्थान नाही असे क्वाइन म्हणतात की,
द्वितीय-क्रम आणि उच्च-क्रम तर्कशास्त्र ही मेंढीचे कातडे पांघरलेल्या संच
उपपत्तीचीच रूपे आहेत; म्हणजेच, विधेयांवर संख्यकीकरण करणाऱ्या प्रणाली पूर्णतः
तार्किक नाहीत.

आत्तासाठी असे तात्त्विक प्रश्न पुढे कधीतरी विचारात घेणेच योग्य ठरेल; आणि
खालील प्रणाली तार्किक असोत वा नसोत, विविध प्रकारच्या युक्तिविचारांची आकारिक
आदर्शीकरणे म्हणून त्यांच्याकडे पाहूया.












\section{बहुवर्गीय तर्कशास्त्र}\label{fol:byd:msl:sec}

प्रथम-क्रम तर्कशास्त्रात चरे आणि संख्यापक एकाच प्रांत वर लागू होतात.
पण अनेक (परस्परविभक्त) प्रांत असणे बऱ्याचदा उपयुक्त ठरते: उदाहरणार्थ,
संख्यांचा प्रांत, भूमितीय वस्तूंचा प्रांत, संख्यांपासून संख्यांकडे
जाणाऱ्या फलनांचा प्रांत, आबेली गटांचा प्रांत, इत्यादी ठेवावेसे वाटू
शकतात.

बहुवर्गीय तर्कशास्त्र अशी चौकट पुरवते. सुरुवात ``वर्गांच्या'' यादीने होते---
वस्तूचा ``वर्ग'' ती कोणत्या ``प्रांत'' मध्ये असणे अपेक्षित आहे ते दर्शवतो.
मग प्रत्येक वर्गासाठी स्वतंत्र चले आणि संख्यापक, तसेच (सामान्यतः)
स्वतंत्र एकरूपता असते. फलने आणि संबंध कोणत्या वर्गांतील वस्तू फलसाधक
म्हणून घेऊ शकतात यानुसार तेही ``प्रकारबद्ध'' केलेले असतात. याखेरीज
प्रथम-क्रम तर्कशास्त्राचे नेहमीचे नियम तसेच ठेवले जातात; संख्यापकांचे नियम
प्रत्येक वर्गासाठी स्वतंत्रपणे पुनरावृत्त होतात.

उदाहरणार्थ, आंतरराष्ट्रीय संबंधांचा अभ्यास करण्यासाठी आपण फ्रेंच नागरिक आणि
जर्मन नागरिक असे वस्तूंचे दोन वर्ग असलेली भाषा निवडू शकतो. पहिल्या वर्गातील
वस्तूंसाठी ``द्राक्षारस पितो'' हा एकस्थानी संबंध; दुसऱ्या वर्गातील वस्तूंसाठी
``वुर्स्ट खातो'' हा दुसरा एकस्थानी संबंध; आणि दोन फलसाधक घेणारा ``बहुराष्ट्रीय
विवाहित जोडपे बनवतात'' हा द्विस्थानी संबंध असू शकतो, ज्याचा पहिला फलसाधक
पहिल्या वर्गातील आणि दुसरा फलसाधक दुसऱ्या वर्गातील असतो. फ्रेंच नागरिकांसाठी
$a$, $b$, $c$ ही चरे आणि जर्मन नागरिकांसाठी $x$, $y$, $z$ ही चरे वापरली, तर
\[\lforall[a][\lforall x][(\Atom{\Obj{MarriedTo}}{a,x} \lif
(\Atom{\Obj{DrinksWine}}{a} \lor \lnot \Atom{\Obj{EatsWurst}}{x})]]
\]
हे सूत्र असे सांगते की कोणत्याही फ्रेंच व्यक्तीचा विवाह एखाद्या जर्मन
व्यक्तीशी झाला असेल, तर एकतर ती फ्रेंच व्यक्ती द्राक्षारस पिते किंवा ती जर्मन
व्यक्ती वुर्स्ट खात नाही.

बहुवर्गीय तर्कशास्त्र प्रथम-क्रम तर्कशास्त्रात स्वाभाविक रीतीने अंतःस्थापित
करता येते: बहुवर्गीय प्रांत मधील सर्व वस्तू एकत्र करून प्रथम-क्रमाचा एक
प्रांत घ्यायचा, वर्गांची नोंद ठेवण्यासाठी एकस्थानी विधेये
वापरायची आणि संख्यापकांना योग्य वर्गापुरते सापेक्ष करायचे. उदाहरणार्थ, वरच्या
उदाहरणाशी संबंधित प्रथम-क्रम भाषेत वर्णिलेल्या इतर संबंधांबरोबर
``$\Obj{German}$'' आणि ``$\Obj{French}$'' ही एकस्थानी विधेये असतील;
मात्र वर्गाविषयीच्या अटी पुसलेल्या असतील. जर्मन वर्गातील चल~$x$
असलेला वर्गबद्ध संख्यापक $\lforall[x][\varphi]$ पुढीलप्रमाणे रूपांतरित होतो:
\[\lforall[x][(\Atom{\Obj{German}}{x} \lif \varphi)].
\]
वर्ग परस्परांपासून वेगळे आहेत याची खात्री देणारी स्वयंसिद्धकेही जोडावी
लागतात---उदाहरणार्थ,
$\lforall[x][\lnot (\Atom{\Obj{German}}{x} \land
  \Atom{\Obj{French}}{x})]$---तसेच ``द्राक्षारस पितो'' हा संबंध फक्त
$\Atom{\Obj{French}}{x}$ या विधेयाची पूर्ति करणाऱ्या वस्तूंनाच लागू होतो
याची हमी देणारी स्वयंसिद्धकेही जोडावी लागतात. या संकेतांमुळे आणि
स्वयंसिद्धकांमुळे बहुवर्गीय वाक्ये प्रथम-क्रम वाक्ये मध्ये आणि
बहुवर्गीय निष्पत्त्या प्रथम-क्रम निष्पत्त्यांमध्ये रूपांतरित होतात,
हे दाखवणे अवघड नाही. तसेच बहुवर्गीय संरचना संबंधित प्रथम-क्रम
संरचनांमध्ये आणि उलट दिशेनेही ``रूपांतरित'' होतात; म्हणून बहुवर्गीय
तर्कशास्त्रासाठीही आपल्याला संपूर्णता प्रमेय मिळते.












\section{द्वितीय-क्रम तर्कशास्त्र}\label{fol:byd:sol:sec}

द्वितीय-क्रम तर्कशास्त्राची भाषा व्यक्तींच्या प्रांत वरच नव्हे, तर त्या
प्रांत वरील संबंधांवरही संख्यापन करण्याची मुभा देते. प्रथम-क्रम
भाषा~$\Lang{L}$ दिली असता, प्रत्येक $k$ साठी $k$-स्थानी संबंधांवर मान घेणारी
$R$ ही चले जोडली जातात आणि त्या चरांवर संख्यापन करण्याची मुभा दिली
जाते. $R$ हे एखाद्या $k$-स्थानी संबंधाचे चल असेल आणि $t_1$,
\dots,~$t_k$ ही नेहमीची (प्रथम-क्रम) पदे असतील, तर
$\Atom{R}{t_1,\dots,t_k}$ हे अण्वीय सूत्र आहे. याखेरीज सूत्रांचा
संच प्रथम-क्रम तर्कशास्त्राप्रमाणेच, फक्त द्वितीय-क्रम संख्यापनासाठी अतिरिक्त
उपनियम घालून, परिभाषित केला जातो. लक्षात घ्या की आपल्याकडे एकरूपता फक्त
प्रथम-क्रम पदांसाठी आहे: $R$ आणि $S$ ही समान स्थानसंख्या~$k$ असलेल्या
संबंधांची चले असतील, तर $\eq[R][S]$ हे पुढील सूत्राचे संक्षिप्त रूप
म्हणून परिभाषित करता येते:
\[\lforall[x_1 \dots][\lforall[x_k][(\Atom{R}{x_1, \dots, x_k} \liff
  \Atom{S}{x_1, \dots, x_k})]].
\]

द्वितीय-क्रम तर्कशास्त्राचे नियम संख्यापकांचे नियम नव्या द्वितीय-क्रम चरांपर्यंत
विस्तारित करतात. मात्र ही चरे $\Lang{L}$ मधील विधेये शी आणि अधिक
सामान्यपणे $\Lang{L}$ मधील सूत्रे शी कशी परस्परक्रिया करतात, हे येथे
थोडे काळजीपूर्वक स्पष्ट करावे लागते. किमान चौकटीत संबंधचरे पदे म्हणून गणली
जातात; म्हणून पुढील प्रकारच्या अनुमानांना मुभा असते:
\[\varphi(R) \Proves \lexists[R][\varphi(R)]
\]
पण $\Lang{L}$ ही अंकगणिताची भाषा असून $<$ हे स्थिर संबंधचिन्ह असेल, तर पुढील
अनुमानही ग्राह्य असावे अशी अपेक्षा असते:
\[x < y \Proves \lexists[R][\Atom{R}{x,y}]
\]
किंवा दिलेल्या सूत्र~$\varphi$ साठी,
\[\varphi(x_1, \dots, x_k) \Proves \lexists[R][\Atom{R}{x_1,\dots,x_k}]
\]
आणखी सामान्यपणे, पुढील प्रकारच्या अनुमानांना मुभा द्यावीशी वाटू शकते:
\[\Subst{\varphi}{\lambd[\vec x][\psi(\vec x)]}{R} \Proves
\lexists[R][\varphi]
\]
येथे $\Subst{\varphi}{\lambd[\vec x][\psi(\vec x)]}{R}$ म्हणजे $\varphi$ मधील
$\Obj{R}{t_1,\dots,t_k}$ या रूपाचे प्रत्येक अण्वीय सूत्र
$\psi(t_1, \dots, t_k)$ ने बदलल्यावर मिळणारे सूत्र. हा शेवटचा नियम
{\em गुणधर्माधारित संबंधरचना-रूपबंध} असण्याशी समतुल्य आहे; म्हणजे
\[\lexists[R][\lforall[x_1, \dots, x_k][(\varphi(x_1, \dots, x_k) \liff
\Atom{R}{x_1, \dots, x_k})]],
\]
या रूपाचे, द्वितीय-क्रम भाषेतील प्रत्येक सूत्र~$\varphi$ साठी एक,
स्वयंसिद्धक असते; त्यात $R$ हे मुक्त चर नसते. (सराव: $R$ ला $\varphi$ मध्ये येऊ
दिल्यास हा रूपबंध विसंगत होतो, हे दाखवा!)

तर्कशास्त्रज्ञ “द्वितीय-क्रम तर्कशास्त्राची स्वयंसिद्धके” म्हणतात तेव्हा
सहसा प्रथम-क्रम तर्कशास्त्राचा द्वितीय-क्रम संख्यापक-नियमांनी केलेला किमान
विस्तार आणि गुणधर्माधारित संबंधरचना-रूपबंध त्यांना अभिप्रेत असतो. पण या
स्वयंसिद्धकांच्या आणि नियमांच्या दुर्बल उपप्रणालींचा अभ्यास करणेही अनेकदा
उपयुक्त असते. उदाहरणार्थ, गुणधर्माधारित संबंधरचनेचा स्वयंसिद्धक-रूपबंध पूर्ण
सामान्यतेत \emph{स्वसमावेशक} आहे: द्वितीय-क्रम संख्यापक असलेल्या
सूत्र ने “परिभाषित” होणारा $\Atom{R}{x_1, \dots, x_k}$ हा संबंध
अस्तित्वात आहे असे विधान करण्याची तो मुभा देतो; आणि हे संख्यापक अशा सर्व
संबंधांच्या संचावर मान घेतात---ज्या संचात $R$ स्वतःही येतो!{} विसाव्या
शतकाच्या सुमारास रसेलच्या विरोधापत्तीला दिलेली एक सामान्य प्रतिक्रिया म्हणजे
अशा परिभाषांनाच तिच्यासाठी दोषी धरणे आणि गणिताचा पाया विकसित करताना त्या
टाळणे. सूत्र~$\varphi$ मध्ये द्वितीय-क्रम संख्यापक वापरण्यास मनाई केली, तर
गुणधर्माधारित संबंधरचनेचे काहीसे दुर्बल असे {\em स्वसमावेशरहित} रूप मिळते.

चिन्हार्थाच्या दृष्टीने द्वितीय-क्रम संरचनांमध्ये त्या भाषेसाठीची एक
प्रथम-क्रम संरचना आणि ज्या संबंधांवर द्वितीय-क्रम संख्यापक मान घेतात
असा त्या प्रांत वरील संबंधांचा संच असतो (अधिक नेमके सांगायचे, तर प्रत्येक
$k$ साठी स्थानसंख्या~$k$ असलेल्या संबंधांचा एक संच असतो). अर्थात,
निष्पत्ती प्रणालीत गुणधर्माधारित संबंधरचना समाविष्ट असेल, तर
“द्वितीय-क्रम भागात” ती स्वयंसिद्धके पूर्ण करण्याइतके संबंध असावेत ही
अतिरिक्त अट येते---नाहीतर निष्पत्ती प्रणाली निर्दोष राहत नाही!{} पुरेसे
संबंध उपलब्ध असल्याची खात्री करण्याचा एक सोपा मार्ग म्हणजे द्वितीय-क्रम भागात
प्रथम-क्रम भागावरील \emph{सर्व} संबंध घेणे. अशा संरचना ला
\emph{पूर्ण} म्हणतात आणि एका अर्थाने तीच भाषेची “अभिप्रेत संरचना”
असते. आपण फक्त पूर्ण संरचनेचा विचार केला, तर त्याला
\emph{पूर्ण} द्वितीय-क्रम चिन्हार्थमीमांसा म्हणतात. अशा वेळी संरचना
निर्दिष्ट करणे म्हणजे फक्त तिचा प्रथम-क्रम भाग निर्दिष्ट करणे; कारण
द्वितीय-क्रम भागातील घटक त्यावरून अप्रत्यक्षपणे ठरतात.

सारांशतः, द्वितीय-क्रम तर्कशास्त्राविषयी बोलताना काही संदिग्धता असते.
निष्पत्ती प्रणालीच्या दृष्टीने पुढीलपैकी कोणतीही गोष्ट अभिप्रेत असू शकते:
\begin{enumerate}
\item काही गुणधर्माधारित संबंधरचना-स्वयंसिद्धकांसह “किमान” द्वितीय-क्रम
  निष्पत्ती प्रणाली.
\item पूर्ण गुणधर्माधारित संबंधरचना असलेली “मानक” द्वितीय-क्रम
  निष्पत्ती प्रणाली.
\end{enumerate}
चिन्हार्थाच्या दृष्टीने पुढीलपैकी कशातही स्वारस्य असू शकते:
\begin{enumerate}
\item “दुर्बल” चिन्हार्थमीमांसा, जिच्यात संरचनांमध्ये प्रथम-क्रम
  भाग आणि गुणधर्माधारित संबंधरचना-स्वयंसिद्धके पूर्ण करण्याइतका मोठा
  द्वितीय-क्रम भाग असतो.
\item “मानक” द्वितीय-क्रम चिन्हार्थमीमांसा, जिच्यात फक्त पूर्ण
  संरचनेचा विचार केला जातो.
\end{enumerate}
तर्कशास्त्रज्ञांनी कोणती निष्पत्ती प्रणाली किंवा कोणती चिन्हार्थमीमांसा
अभिप्रेत आहे हे सांगितले नाही, तर सहसा प्रत्येक यादीतील दुसरी बाब त्यांना
अभिप्रेत असते. पुढे दिसेल त्याप्रमाणे, ही चिन्हार्थमीमांसा वापरण्याचा फायदा असा
की तिच्यात अनेक स्वाभाविक गणितीय रचनांची समरूपतेपर्यंत एकमेव वर्णने मिळतात;
त्याच वेळी निष्पत्ती प्रणाली बरीच समर्थ असून या चिन्हार्थमीमांसेसाठी
निर्दोष असते. मात्र निष्पत्ती प्रणाली या चिन्हार्थमीमांसेसाठी
\emph{संपूर्ण नाही}; इतकेच नव्हे, तर प्रभावी रीतीने दिलेली \emph{कोणतीही}
निष्पत्ती प्रणाली पूर्ण द्वितीय-क्रम चिन्हार्थमीमांसेसाठी संपूर्ण नसते.
दुसरीकडे, दुर्बल चिन्हार्थमीमांसेसाठी निष्पत्ती प्रणाली \emph{संपूर्ण
आहे}, हे आपण पाहू; म्हणून एखादे वाक्य सिद्ध करता येत नसेल, तर ते असत्य असणारी
\emph{एखादी तरी} संरचना असते, जरी ती पूर्ण रचना असेलच असे नाही.

द्वितीय-क्रम तर्कशास्त्राची भाषा बरीच अभिव्यक्तिसमर्थ आहे. एकस्थानी संबंधांची
प्रांताच्या उपसंचांशी ओळख करता येते; म्हणून विशेषतः या संचांवर संख्यापनही
करता येते. उदाहरणार्थ, स्वाभाविक संख्यांसाठीचे आगमन एकाच स्वयंसिद्धकाने
व्यक्त करता येते:
\[\lforall[R][((\Atom{R}{\Obj{0}} \land \lforall[x][(\Atom{R}{x} \lif
    \Atom{R}{x'})]) \lif \lforall[x][\Atom{R}{x}])].
\]
अंकगणिताच्या भाषेत $\Obj 0, \Obj \prime, +, \times$ आणि $<$ ही चिन्हे घेतली,
तर त्यांचे वर्तन वर्णिण्यासाठी पुढील स्वयंसिद्धके जोडता येतात:
\begin{enumerate}
\item $\lforall[x][\lnot x' = \Obj 0]$
\item $\lforall[x][\lforall[y][(s(x) = s(y) \lif x = y)]]$
\item $\lforall[x][(x + \Obj 0) = x]$
\item $\lforall[x][\lforall[y][(x + y') = (x + y)']]$
\item $\lforall[x][(x \times \Obj 0) = \Obj 0]$
\item $\lforall[x][\lforall[y][(x \times y') = ((x \times y) + x)]]$
\item $\lforall[x][\lforall[y][(x < y \liff \lexists[z][y = (x + z')])]]$
\end{enumerate}
आपण पूर्ण द्वितीय-क्रम चिन्हार्थमीमांसा वापरत असू, तर ही स्वयंसिद्धके आणि
वरील आगमनाचे स्वयंसिद्धक मिळून अंकगणिताचे मानक प्रतिमान असलेल्या
संरचना~$\Struct{N}$ चे समरूपतेपर्यंत एकमेव वर्णन देतात, हे दाखवणे
अवघड नाही. ही स्वयंसिद्धके सत्य असलेली कोणतीही संरचना~$\Struct{M}$
दिली असता, $\Nat$ वरील नेहमीचे पुनरावर्तन वापरून $\Nat$ पासून
$\Struct{M}$ च्या प्रांत कडे जाणारे $f$ हे फलन असे परिभाषित करा की
$f(0) = \Assign{\Obj 0}{M}$ आणि $f(x+1) = \Assign{\prime}{M}(f(x))$.
$\Nat$ वरील नेहमीचे आगमन आणि $\Struct M$ मध्ये स्वयंसिद्धके (1) आणि~(2)
सत्य असण्याचा उपयोग केल्यावर $f$ हे एकास-एक आहे असे दिसते. $f$ हे
आच्छादक आहे हे पाहण्यासाठी, $f$ च्या व्याप्तीत असलेल्या
$\Domain{M}$ मधील घटकांचा संच $P$ घ्या. $\Struct M$ पूर्ण असल्यामुळे $P$
हा द्वितीय-क्रम प्रांत मध्ये असतो. $f$ च्या रचनेवरून
$\Assign{\Obj 0}{M}$ हा $P$ मध्ये असतो आणि $P$ हा
$\Assign{\prime}{M}$ खाली बंद असतो. $\Struct M$ मध्ये आगमनाचे स्वयंसिद्धक
सत्य असणे (विशेषतः $P$ साठी) $P$ हा $\Struct M$ च्या संपूर्ण प्रथम-क्रम
प्रांत इतकाच असल्याची हमी देते. यावरून $f$ हे एकास-एक आच्छादन आहे.
$\Nat$ वर नेहमीचे आगमन वारंवार वापरून $f$ हे रचनारक्षक प्रतिचित्रण आहे हे
दाखवणे यापेक्षा अवघड नाही.

संचसैद्धान्तिक भाषेत फलन हा संबंधाचाच एक विशेष प्रकार असतो; उदाहरणार्थ,
$\lforall[x][\lexists![y][R(x,y)]]$ पूर्ण करणाऱ्या $R$ या द्विस्थानी
संबंधाशी एकस्थानी फलन~$f$ ची ओळख करता येते. त्यामुळे फलनांवरही संख्यापन करता
येते. मग पूर्ण चिन्हार्थमीमांसा वापरून अनंत संरचनेचा वर्ग असा
परिभाषित करता येतो: त्या संरचना $\Struct M$ असतात ज्यांच्यासाठी
$\Struct M$ च्या प्रांत पासून त्याच्याच उचित उपसंचाकडे जाणारे एक
एकास-एक फलन असते:
\[\lexists[f][(\lforall[x][\lforall[y][(\eq[f(x)][f(y)] \lif \eq[x][y])]]
  \land \lexists[y][\lforall[x][\eq/[f(x)][y]]])].
\]
या वाक्याचा निषेध मग सांत संरचनेचा वर्ग परिभाषित करतो.

याखेरीज, रेषीय क्रमाच्या परिभाषेत पुढील अट जोडून सुक्रमांचा वर्ग परिभाषित
करता येतो:
\[\lforall[P][(\lexists[x][\Atom{P}{x}] \lif \lexists[x][(\Atom{P}{x}
    \land \lforall[y][(y < x \lif \lnot \Atom{P}{y})])])].
\]
“संच” आणि “एकस्थानी संबंध” यांची ओळख केल्यास, प्रत्येक रिक्तेतर संचाला
लघुतम घटक असतो असे हे सांगते. आणखी एका उदाहरणात, शिखरांची परस्पर न जोडलेल्या
भागांत कोणतीही अ-तुच्छ विभागणी नसते असे सांगून आलेखांची संयोजकता व्यक्त करता
येते:
\[\lnot \lexists[A][(\lexists[x][A(x)] \land \lexists[y][\lnot A(y)]
  \land \lforall[w][\lforall[z][((\Atom{A}{w} \land \lnot \Atom{A}{z})
      \lif \lnot \Atom{R}{w,z})]])].
\]
दुसऱ्या उदाहरणासाठी, ज्या सांत संरचनेच्या प्रांताचा आकार सम
आहे त्यांचा वर्ग परिभाषित करण्याचा सराव तुम्ही करू शकता. आणखी लक्षवेधक
म्हणजे, परिमेय संख्या अंतर्भूत करणारे संपूर्ण क्रमित क्षेत्र म्हणून वास्तव
संख्यांचे समरूपतेपर्यंत एकमेव वर्णन देता येते.

थोडक्यात, द्वितीय-क्रम तर्कशास्त्र प्रथम-क्रम तर्कशास्त्रापेक्षा खूपच अधिक
अभिव्यक्तिसमर्थ आहे. ही चांगली बाजू; आता अडचण पाहू. पूर्ण द्वितीय-क्रम
चिन्हार्थमीमांसेसाठी संपूर्ण असणारी प्रभावी निष्पत्ती प्रणाली नाही, हे
आपण आधीच नमूद केले आहे. प्रथम-क्रम तर्कशास्त्राचे अनेक गुणधर्म येथे उपलब्ध
नसतात; त्यांत संहतता आणि लोव्हेनहाइम--स्कोलेम प्रमेयांचा समावेश होतो.

दुसरीकडे, पूर्ण द्वितीय-क्रम चिन्हार्थमीमांसा सोडून दुर्बल चिन्हार्थमीमांसा
स्वीकारली, तर किमान द्वितीय-क्रम निष्पत्ती प्रणाली या
चिन्हार्थमीमांसेसाठी संपूर्ण असते. दुसऱ्या शब्दांत, $\Proves$ चा अर्थ
“किमान प्रणालीत सिद्ध होते” आणि $\Entails$ चा अर्थ “दुर्बल
चिन्हार्थमीमांसेत तार्किकरीत्या निष्पन्न होते” असा घेतल्यास, $\Gamma \Entails
\varphi$ असेल तेव्हा $\Gamma \Proves \varphi$ असते, हे दाखवता येते.
निष्पत्ती प्रणालीत ठरावीक गुणधर्माधारित संबंधरचना-स्वयंसिद्धके समाविष्ट
करायची असतील, तर ही स्वयंसिद्धके पूर्ण करणाऱ्या द्वितीय-क्रम रचनांपुरती
चिन्हार्थमीमांसा मर्यादित करावी लागते: उदाहरणार्थ, $\Delta$ हा
गुणधर्माधारित संबंधरचना-स्वयंसिद्धकांचा संच (शक्यतो त्यांपैकी सर्वांचा संच)
असेल, तर $\Gamma \cup \Delta \Entails \varphi$ असल्यास $\Gamma \cup \Delta
\Proves \varphi$ असते. विशेषतः, आपण विचारात घेत असलेली गुणधर्माधारित
संबंधरचना-स्वयंसिद्धके वापरून $\varphi$ सिद्ध करता येत नसेल, तर ती स्वयंसिद्धके
सत्य असूनही $\lnot \varphi$ सत्य असणारे एक प्रतिमान असते.

दुर्बल चिन्हार्थमीमांसेसाठी संपूर्णता प्रमेय का लागू होते हे पाहण्याचा सर्वांत
सोपा मार्ग म्हणजे द्वितीय-क्रम तर्कशास्त्राचा पुढीलप्रमाणे बहुवर्गीय
तर्कशास्त्र म्हणून विचार करणे. एका वर्गाचा अर्थ नेहमीचा “प्रथम-क्रम”
प्रांत असा लावला जातो; मग प्रत्येक $k$ साठी “स्थानसंख्या $k$ असलेल्या
संबंधांचा” प्रांत असतो. भाषेत
“$\Atom{\Obj{true}_k}{R,x_1,\dots,x_k}$” ही अंगभूत संबंधचिन्हे घेतो. येथे
$R$ हे “$k$-स्थानी संबंध” या वर्गाचे चर आणि $x_1$, \dots,~$x_k$ या
प्रथम-क्रम वर्गातील वस्तू असताना, $R$ हा संबंध $x_1$, \dots,~$x_k$ यांना
लागू होतो असे हे चिन्ह सांगते.

ही ओळख स्वीकारल्यावर दुर्बल द्वितीय-क्रम चिन्हार्थमीमांसा मूलतः बहुवर्गीय
तर्कशास्त्राच्या नेहमीच्या चिन्हार्थमीमांसेसारखीच असते; आणि बहुवर्गीय
तर्कशास्त्राचे प्रथम-क्रम तर्कशास्त्रात अंतःस्थापन करता येते, हे आपण आधीच
पाहिले आहे. त्यामुळे दोन्ही दिशांतील रूपांतरणांचा विचार करता, द्वितीय-क्रम
तर्कशास्त्राची दुर्बल संकल्पना हे प्रत्यक्षात प्रथम-क्रम तर्कशास्त्राचेच
छद्मरूप आहे; त्यात प्रांत मध्ये योग्य स्वयंसिद्धकांच्या अधीन असलेल्या
“वस्तू” आणि “संबंध” या दोन्हींचा समावेश असतो.












\section{उच्च-क्रम तर्कशास्त्र}\label{fol:byd:hol:sec}

प्रथम-क्रम तर्कशास्त्रातून द्वितीय-क्रम तर्कशास्त्राकडे गेल्यामुळे औपचारिक
भाषेच्या चौकटीत प्रथम-क्रम प्रांत मधील वस्तूंच्या संचांविषयी बोलणे शक्य
झाले. तेथेच का थांबावे? उदाहरणार्थ, तृतीय-क्रम तर्कशास्त्रामुळे वस्तूंच्या
संचांचे संच, किंवा वस्तू आणि वस्तूंचे संच हे दोन्ही अंतर्भूत करणारे संचही
हाताळता यावेत. चतुर्थ-क्रम तर्कशास्त्रामुळे अशा प्रकारच्या वस्तूंच्या
संचांविषयी बोलता येईल. तुमच्या लक्षात आले असेलच की ही कल्पना कितीही वेळा
पुनरावृत्त करता येते.

व्यवहारात उच्च-क्रम तर्कशास्त्र संबंधांऐवजी फलनांच्या मदतीने
सूत्रted केले जाते. (स्वाभाविक ओळखी गृहीत धरल्यास हा फरक महत्त्वाचा
नाही.) काही मूलभूत “वर्ग” $A$, $B$, $C$,~\dots (ज्यांना आता आपण “प्रकार”
म्हणू) दिले असता, पुढील अट घालून नवे प्रकार तयार करता येतात:
\begin{quote}
$\sigma$ आणि $\tau$ हे सांत प्रकार असतील, तर $\sigma \to \tau$ हाही सांत
प्रकार असतो.
\end{quote}
प्रकारांना विन्यासात्मक “लेबले” समजा; आपल्या प्रांत मध्ये हव्या असलेल्या
वस्तूंचे ती वर्गीकरण करतात. $\sigma \to \tau$ हा प्रकार, प्रकार~$\sigma$ च्या
वस्तू घेऊन प्रकार~$\tau$ च्या वस्तू देणारी फलने असलेल्या वस्तूंचे वर्णन करतो.
उदाहरणार्थ, “सत्य” आणि “असत्य” या सत्यतामूल्यांचा $\Omega$ हा प्रकार आणि
नैसर्गिक संख्यांचा $\Nat$ हा प्रकार घ्यावासा वाटू शकतो. अशा वेळी $\Nat \to
\Omega$ प्रकारच्या वस्तू एकस्थानी संबंध किंवा $\Nat$ चे उपसंच मानता येतात;
$\Nat \to \Nat$ प्रकारच्या वस्तू ही नैसर्गिक संख्यांपासून नैसर्गिक संख्यांकडे
जाणारी फलने असतात; आणि $(\Nat \to \Nat) \to \Nat$ प्रकारच्या वस्तू
“फलनके” असतात, म्हणजे फलने घेऊन संख्या देणारी उच्च-प्रकार फलने.

द्वितीय-क्रम तर्कशास्त्राप्रमाणेच, उच्च-क्रम तर्कशास्त्राकडे बहुवर्गीय
तर्कशास्त्राचा एक प्रकार म्हणून पाहता येते; त्यात विचारात घ्यायच्या प्रत्येक
वस्तुप्रकारासाठी एक वर्ग असतो. पण उच्च-प्रकार तर्कशास्त्राचा विन्यास
मुळापासूनच परिभाषित करणे सहसा अधिक स्पष्ट ठरते. उदाहरणार्थ, सांत प्रकारांचा
संच पुढीलप्रमाणे विगमनाने परिभाषित करता येतो:
\begin{enumerate}
\item $\Nat$ हा सांत प्रकार आहे.
\item $\sigma$ आणि $\tau$ हे सांत प्रकार असतील, तर $\sigma
  \to \tau$ हाही सांत प्रकार आहे.
\item $\sigma$ आणि $\tau$ हे सांत प्रकार असतील, तर $\sigma \times
  \tau$ हाही सांत प्रकार आहे.
\end{enumerate}
सहज अर्थाने, $\Nat$ नैसर्गिक संख्यांचा प्रकार दर्शवतो, $\sigma \to \tau$
हा $\sigma$ पासून $\tau$ कडे जाणाऱ्या फलनांचा प्रकार दर्शवतो आणि
$\sigma \times \tau$ हा वस्तूंच्या जोड्यांचा प्रकार दर्शवतो, ज्यातील एक वस्तू
$\sigma$ मधील आणि दुसरी $\tau$ मधील असते. मग पदांचा संच पुढीलप्रमाणे
विगमनाने परिभाषित करता येतो:
\begin{enumerate}
\item प्रत्येक प्रकार $\sigma$ साठी $x$, $y$, $z$, \dots ही प्रकार~$\sigma$
  ची अनेक चरे असतात.
\item $\Obj 0$ हे $\Nat$ प्रकारचे पद आहे.
\item $\Obj S$ (उत्तरवर्ती) हे $\Nat \to \Nat$ प्रकारचे पद आहे.
\item $s$ हे $\sigma$ प्रकारचे पद आणि $t$ हे
  $\Nat \to (\sigma \to \sigma)$ प्रकारचे पद असेल, तर $\Obj{R}_{st}$ हे
  $\Nat \to \sigma$ प्रकारचे पद आहे.
\item $s$ हे $\tau \to \sigma$ प्रकारचे पद आणि $t$ हे प्रकार~$\tau$ चे पद
  असेल, तर $s(t)$ हे $\sigma$ प्रकारचे पद आहे.
\item $s$ हे प्रकार~$\sigma$ चे पद आणि $x$ हे प्रकार~$\tau$ चे चर असेल,
  तर $\lambd[x][s]$ हे $\tau \to \sigma$ प्रकारचे पद आहे.
\item $s$ हे प्रकार~$\sigma$ चे पद आणि $t$ हे प्रकार~$\tau$ चे पद असेल,
  तर $\tuple{s, t}$ हे $\sigma \times \tau$ प्रकारचे पद आहे.
\item $s$ हे प्रकार~$\sigma \times \tau$ चे पद असेल, तर $p_1(s)$ हे
  प्रकार~$\sigma$ चे पद आणि $p_2(s)$ हे प्रकार~$\tau$ चे पद आहे.
\end{enumerate}
सहज अर्थाने, $\Obj{R}_{st}$ हे पुढीलप्रमाणे पुनरावर्ती रीतीने परिभाषित केलेले
फलन दर्शवते:
\begin{align*}
\Obj{R}_{st}(0) & = s \\
\Obj{R}_{st}(x+1) & = t(x, R_{st}(x)),
\end{align*}
$\tuple{s, t}$ ही ज्याचा पहिला घटक~$s$ आणि दुसरा घटक~$t$ आहे अशी जोडी
दर्शवते, तर $p_1(s)$ आणि~$p_2(s)$ हे $s$ चे पहिले व दुसरे घटक
(“प्रक्षेपणे”) दर्शवतात. शेवटी, $\lambd[x][s]$ हे
\[f(x) = s
\]
या अटीने प्रकार~$\tau$ च्या कोणत्याही~$x$ साठी परिभाषित झालेले फलन~$f$
दर्शवते.

म्हणून बाब (6) आपल्याला गुणधर्माधारित फलनरचनेचे एक रूप देते आणि
पदांच्या मदतीने फलने परिभाषित करणे शक्य करते. समान प्रकारच्या पदांमधील
एकरूपता विधाने $\eq[s][t]$, नेहमीचे विधानीय संयोजक आणि उच्च-प्रकार
संख्यापन यांपासून सूत्रे घडवली जातात. वर परिभाषित केलेल्या पदांचे
नियमन करणारी मूलभूत समीकरणे, तसेच संख्यापक आणि एकरूपता असलेले
तर्कशास्त्राचे नेहमीचे नियम, ही मग प्रणालीची स्वयंसिद्धके मानता येतात.

सांत प्रकारांच्या प्रणालीत सत्यतामूल्यांचा $\Omega$ हा प्रकार जोडला, तर
त्याचा वापर नियंत्रित करणारी स्वयंसिद्धकेही समाविष्ट करावी लागतात. खरे तर
योग्य युक्ती वापरून गुंतागुंतीची सूत्रे पूर्णपणे काढून टाकता येतात आणि
त्यांच्या जागी $\Omega$ प्रकारची पदे ठेवता येतात!{} त्यानुसार सिद्धता-पद्धतीतही
बदल करता येतो. त्यातून मूलतः १९३० च्या दशकात ॲलोन्झो चर्च यांनी मांडलेली
\emph{प्रकारांची साधी उपपत्ती} मिळते.

द्वितीय-क्रम तर्कशास्त्राप्रमाणेच, उच्च-प्रकार चिन्हार्थमीमांसेच्याही विविध
आवृत्त्या वापरता येतात. पूर्ण आवृत्तीत $\sigma \to \tau$ प्रकारची चरे,
प्रकार~$\sigma$ च्या वस्तूंपासून प्रकार~$\tau$ च्या वस्तूंकडे जाणाऱ्या
\emph{सर्व} फलनांच्या संचावर मान घेतात. अपेक्षेप्रमाणे ही चिन्हार्थमीमांसा
इतकी समर्थ आहे की तिच्यासाठी संपूर्ण आणि प्रभावी निष्पत्ती प्रणाली
असू शकत नाही. पण दुर्बल चिन्हार्थमीमांसेचाही विचार करता येतो; त्यात
संरचनांमध्ये प्रत्येक प्रकार $\tau$ साठी $T_\tau$ हे घटकांचे संच
आणि उपयोजन, प्रक्षेपण इत्यादींसाठी योग्य कृत्ये असतात. तपशील योग्य रीतीने
मांडल्यास वर वर्णिलेल्या प्रकारच्या निष्पत्ती प्रणालींसाठी संपूर्णता
प्रमेये मिळवता येतात.

उच्च-प्रकार तर्कशास्त्र आकर्षक आहे, कारण गणिताचा मोठा भाग स्वाभाविक रीतीने
अंतःस्थापित करता येईल अशी चौकट ते देते: $\Nat$ पासून सुरुवात करून वास्तव
संख्या, संतत फलने इत्यादी परिभाषित करता येतात. अंतःप्रज्ञावादी तर्कशास्त्राच्या
संदर्भातही ते विशेष आकर्षक आहे, कारण प्रकारांना स्पष्ट “रचनाशील” अर्थनिर्धारणे
आहेत. खरे तर, उच्च-प्रकार चिन्हार्थमीमांसेच्या रचनाशील आवृत्त्या
(अभिजात तर्कशास्त्राऐवजी अंतःप्रज्ञावादी तर्कशास्त्रावर आधारलेल्या) विकसित
करता येतात. त्या ही रचनाशील अर्थनिर्धारणे अतिशय स्पष्ट करतात आणि अनेक
दृष्टींनी त्यांच्या अभिजात समकक्षांपेक्षा अधिक रोचक असतात.












\section{अंतःप्रज्ञावादी तर्कशास्त्र}\label{fol:byd:il:sec}

द्वितीय-क्रम आणि उच्च-क्रम तर्कशास्त्राच्या उलट, अंतःप्रज्ञावादी प्रथम-क्रम
तर्कशास्त्र हे अभिजात आवृत्तीचे एक निर्बंधित रूप आहे; अधिक “रचनाशील” प्रकारच्या
तर्कविचाराचे प्रतिमान देणे हा त्याचा उद्देश आहे. त्यामागील काही प्रेरणा पुढील
उदाहरणांनी स्पष्ट होतील.

समजा, एखाद्या दिवशी कोणीतरी तुमच्याकडे येऊन सांगितले की त्यांनी अशी नैसर्गिक
संख्या~$x$ ठरवली आहे की $x$ अविभाज्य असेल तर रीमान गृहीतक सत्य असते आणि $x$
संयुक्त असेल तर रीमान गृहीतक असत्य असते. फारच चांगली बातमी!{} रीमान गृहीतक
सत्य आहे की नाही हा गणितातील मोठ्या अनिर्णित प्रश्नांपैकी एक आहे; आणि येथे तर
ही समस्या केवळ गणनेच्या समस्येत, म्हणजे एखादी ठरावीक संख्या अविभाज्य आहे की
नाही हे ठरवण्यात, रूपांतरित झाल्यासारखी दिसते.

$x$ चे ते अद्भुत मूल्य कोणते? ते त्याचे पुढीलप्रमाणे वर्णन करतात: रीमान गृहीतक
सत्य असेल तर $x$ ही $7$ इतकी आणि अन्यथा $9$ इतकी नैसर्गिक संख्या आहे.

तुम्ही चिडून आपले पैसे परत मागता. अभिजात दृष्टिकोनातून वरील वर्णन $x$ चे
एकमेव मूल्य खरोखर ठरवते; पण तुम्हाला प्रत्यक्षात हवे आहे ते
\emph{स्पष्टपणे} दिलेले $x$ चे मूल्य.

आणखी एक, कदाचित कमी कृत्रिम, उदाहरण पाहू. अपरिमेय संख्येचा परिमेय घात घेऊन
परिमेय उत्तर मिळू शकते, हे आपल्याला माहीत आहे. उदाहरणार्थ,
$\sqrt{2}^2 = 2$. अपरिमेय संख्येचा \emph{अपरिमेय} घात घेऊन परिमेय उत्तर
मिळू शकते की नाही, हे तितके स्पष्ट नाही. पुढील प्रमेय याचे होकारार्थी उत्तर
देते:

\begin{thm}
$a^b$ परिमेय असेल अशा $a$ आणि $b$ या अपरिमेय संख्या आहेत.
\end{thm}

\begin{proof}
$\sqrt{2}^{\sqrt{2}}$ चा विचार करा. ही संख्या परिमेय असेल, तर आपले काम झाले:
$a = b = \sqrt{2}$ घेता येईल. अन्यथा ती अपरिमेय आहे. मग
\[(\sqrt{2}^{\sqrt{2}})^{\sqrt{2}} = \sqrt{2}^{\sqrt{2} \cdot
  \sqrt{2}} = \sqrt{2}^2 = 2,
\]
आणि ही संख्या निश्चितच परिमेय आहे. म्हणून या प्रकरणात $a$ म्हणून
$\sqrt{2}^{\sqrt{2}}$ आणि $b$ म्हणून~$\sqrt 2$ घ्या.
\end{proof}

याला वैध सिद्धता म्हणता येईल का? बहुतेक गणितज्ञांच्या मते होय. पण यात पुन्हा
काहीसे असमाधानकारक असे काही आहे: विशिष्ट गुणधर्म असलेल्या वास्तव संख्यांच्या
जोडीचे अस्तित्व आपण सिद्ध केले, पण ती संख्यांची \emph{कोणती} जोडी आहे हे
सांगता आलेले नाही. हाच निष्कर्ष अशा रीतीने सिद्ध करता येतो की सिद्धतेत
$a$, $b$ ही जोडी \emph{प्रत्यक्ष दिलेली} असते: $a = \sqrt{3}$ आणि
$b = \log_3 4$ घ्या. मग
\[a^b = \sqrt{3}^{\log_3 4} = 3^{1/2 \cdot \log_3 4} = (3^{\log_3
  4})^{1/2} = 4^{1/2}= 2,
\]
कारण $3^{\log_3 x} = x$.

पहिल्या सिद्धतेतील पाऊल अमान्य असणाऱ्या तर्कविचाराचे प्रतिमान देण्यासाठी
अंतःप्रज्ञावादी तर्कशास्त्र रचले आहे. $\varphi(x)$ ची पूर्ति करणारा $x$ अस्तित्वात
आहे हे सिद्ध करणे म्हणजे दुसऱ्या सिद्धतेप्रमाणे एखादा ठरावीक~$x$ आणि तो
$\varphi$ ची पूर्ति करतो याची सिद्धता द्यावी लागते. $\varphi$ किंवा $\psi$ सत्य आहे हे
सिद्ध करण्यासाठी त्यांपैकी एकाची सिद्धता देता आली पाहिजे.

आकारिक भाषेत सांगायचे, तर प्रथम-क्रम तर्कशास्त्राची निष्पत्ती प्रणाली
ठरावीक रीतीने निर्बंधित केल्यावर अंतःप्रज्ञावादी प्रथम-क्रम तर्कशास्त्र मिळते.
त्याचप्रमाणे द्वितीय-क्रम किंवा उच्च-क्रम तर्कशास्त्राच्याही अंतःप्रज्ञावादी
आवृत्त्या असतात. गणितीय दृष्टिकोनातून या केवळ आकारिक निगमन-पद्धती आहेत; पण
आधी नमूद केल्याप्रमाणे, एका प्रकारच्या गणितीय तर्कविचाराचे प्रतिमान देणे हा
त्यांचा उद्देश आहे. हा तर्कविचार गणिताविषयीच्या एखाद्या तात्त्विक दृष्टिकोनाने
(उदाहरणार्थ, ब्रौवर यांच्या अंतःप्रज्ञावादाने) समर्थित मानता येतो; किंवा
बिशप यांच्या रचनाशीलतावादाच्या धर्तीवर अधिक “मूर्त” आणि समाधानकारक असा
गणितीय तर्कविचार मानता येतो; तसेच आकारिक वर्णन अनौपचारिक प्रेरणा खरोखर
पकडते की नाही यावर वादही घालता येतो. पण आपली तात्त्विक भूमिका कोणतीही असो,
अंतःप्रज्ञावादी तर्कशास्त्राचा आकारिक रीतीने मांडलेले तर्कशास्त्र म्हणून
अभ्यास करता येतो; आणि कारणे कोणतीही असोत, अनेक गणितीय तर्कशास्त्रज्ञांना तो
अभ्यास रंजक वाटतो.

अंतःप्रज्ञावादी संयोजकांचे एक अनौपचारिक रचनाशील अर्थनिर्धारण आहे. ते सहसा
BHK अर्थनिर्धारण म्हणून ओळखले जाते (ब्रौवर, हीटिंग आणि कोल्मोगोरोव यांच्या
नावांवरून). ते असे आहे: $\varphi \land \psi$ ची सिद्धता म्हणजे $\varphi$ ची सिद्धता
आणि $\psi$ ची सिद्धता यांची जोडी; $\varphi \lor \psi$ ची सिद्धता म्हणजे $\varphi$ ची
सिद्धता किंवा $\psi$ ची सिद्धता, तसेच यांपैकी नेमके कोणते प्रकरण आहे याची स्पष्ट
माहिती; $\varphi \lif \psi$ ची सिद्धता म्हणजे $\varphi$ च्या सिद्धतेचे $\psi$ च्या
सिद्धतेत रूपांतर करणारी प्रक्रिया; $\lforall[x][\varphi(x)]$ ची सिद्धता म्हणजे
$x$ च्या कोणत्याही मूल्यासाठी $\varphi(x)$ ची सिद्धता देणारी प्रक्रिया; आणि
$\lexists[x][\varphi(x)]$ ची सिद्धता म्हणजे $x$ चे एखादे मूल्य आणि ते मूल्य $\varphi$
ची पूर्ति करते याची सिद्धता. या अर्थनिर्धारणाचे संगणकीय वर्णन
“करी--हॉवर्ड समरूपण” किंवा “सूत्रे-हे-प्रकार प्रतिमान” या नावांनी करता
येते: सूत्र एखादा विशिष्ट दत्तांश-प्रकार ठरवते असे समजा आणि
सिद्धतांना त्या दत्तांश-प्रकारांच्या संगणकीय वस्तू समजा; या वस्तूंमुळे संबंधित
सूत्र सत्य आहे हे आपल्याला दिसते.

अंतःप्रज्ञावादी तर्कशास्त्राकडे अनेकदा “विमध्य सिद्धांत वजा केल्यावर उरणारे”
अभिजात तर्कशास्त्र म्हणून पाहिले जाते. पुढील प्रमेय हे अधिक नेमके करते.
\begin{thm}
अंतःप्रज्ञावादी तर्कशास्त्रात पुढील स्वयंसिद्धक-रूपबंध समतुल्य आहेत:
\begin{enumerate}
\item $(\lnot \varphi \lif \lfalse) \lif \varphi$.
\item $\varphi \lor \lnot \varphi$
\item $\lnot \lnot \varphi \lif \varphi$
\end{enumerate}
\end{thm}
कोणत्याही एका रूपबंधापासून इतर दोन्हींचे दृष्टांत मिळवणे हा अंतःप्रज्ञावादी
तर्कशास्त्रातील चांगला सराव आहे.

ब्रौवर यांच्या अंतःप्रज्ञावादाचे आकारीकरण म्हणून मांडलेल्या अंतःप्रज्ञावादी
विधानीय तर्कशास्त्राच्या पहिल्या निगमन-पद्धती कोल्मोगोरोव, ग्लिवेन्को आणि
हीटिंग यांनी स्वतंत्रपणे दिल्या. अंतःप्रज्ञावादी प्रथम-क्रम तर्कशास्त्राचे
(आणि अंतःप्रज्ञावादी गणिताच्या काही भागांचे) पहिले आकारीकरण हीटिंग यांनी
केले. अभिजात तर्कशास्त्रात वैध असलेले अनेक रूपबंध अंतःप्रज्ञावादी
तर्कशास्त्रात वैध नसले, तरी त्यांपैकी अनेक वैधही असतात.

\emph{द्विनकरण रूपांतरण} अभिजात आणि अंतःप्रज्ञावादी तर्कशास्त्रांतील एक
महत्त्वाचा संबंध वर्णिते. ते पुढीलप्रमाणे विगमनाने परिभाषित केले जाते
($\varphi^N$ हे अभिजात सूत्र~$\varphi$ चे “अंतःप्रज्ञावादी” रूपांतर समजा):
\begin{align*}
\varphi^N & \ident \lnot\lnot \varphi \quad \text{अण्वीय सूत्रे $\varphi$ साठी} \\
(\varphi \land \psi)^N & \ident (\varphi^N \land \psi^N) \\
(\varphi \lor \psi)^N & \ident  \lnot\lnot (\varphi^N \lor \psi^N) \\
(\varphi \lif \psi)^N & \ident (\varphi^N \lif \psi^N) \\
(\lforall[x][\varphi])^N & \ident \lforall[x][\varphi^N] \\
(\lexists[x][\varphi])^N & \ident \lnot\lnot\lexists[x][\varphi^N]
\end{align*}
कोल्मोगोरोव आणि ग्लिवेन्को यांनी विधानीय तर्कशास्त्रासाठी या रूपांतरणाच्या
आवृत्त्या दिल्या होत्या; विधेय तर्कशास्त्रासाठी ते गोडेल आणि गेंटझेन यांनी
स्वतंत्रपणे दिले. आपल्याला पुढील निष्कर्ष मिळतो.

\begin{thm}
\begin{enumerate}
\item $\varphi \liff \varphi^N$ हे अभिजात तर्कशास्त्रात सिद्धतायोग्य आहे.
\item $\varphi$ हे अभिजात तर्कशास्त्रात सिद्धतायोग्य असेल, तर $\varphi^N$ हे
  अंतःप्रज्ञावादी तर्कशास्त्रात सिद्धतायोग्य आहे.
\end{enumerate}
\end{thm}

आता पुढील संवादाची कल्पना करता येते. अभिजात गणितज्ञ: “मी $\varphi$ सिद्ध केले
आहे!” अंतःप्रज्ञावादी गणितज्ञ: “तुमची सिद्धता वैध नाही. तुम्ही प्रत्यक्षात
$\varphi^N$ सिद्ध केले आहे.” अभिजात गणितज्ञ: “मला तेही चालेल!” अभिजात
गणितज्ञाच्या दृष्टीने अंतःप्रज्ञावादी गणितज्ञ क्षुल्लक भेद करत आहे, कारण ती
दोन्ही समतुल्य आहेत. पण अंतःप्रज्ञावादी गणितज्ञाच्या मते त्यांत फरक आहे.

लक्षात घ्या की वरील रूपांतरण फक्त शुद्ध तर्कशास्त्राशी संबंधित आहे; अभिजात
आणि अंतःप्रज्ञावादी गणितासाठी योग्य \emph{अतार्किक} स्वयंसिद्धके कोणती किंवा
त्यांच्यात कोणता संबंध आहे, या प्रश्नाला ते हात घालत नाही. पण वरील प्रमेयाचा
पुढील किंचित विस्तार काही उपयुक्त माहिती देतो:

\begin{thm}
$\Gamma$ ने $\varphi$ अभिजात रीतीने सिद्ध होत असेल, तर $\Gamma^N$ ने $\varphi^N$
अंतःप्रज्ञावादी रीतीने सिद्ध होते.
\end{thm}

दुसऱ्या शब्दांत, काही गृहीतकांपासून $\varphi$ अभिजात रीतीने सिद्धतायोग्य असेल,
तर त्यांच्या द्विनकरण रूपांतरांपासून $\varphi^N$ अंतःप्रज्ञावादी रीतीने
सिद्धतायोग्य असते.

एखादे वाक्य किंवा विधानीय सूत्र अंतःप्रज्ञावादी रीतीने वैध आहे हे
दाखवण्यासाठी फक्त त्याची सिद्धता द्यावी लागते. पण ते वैध नाही हे कसे
दाखवणार? त्यासाठी आपल्याला निर्दोष, आणि शक्यतो संपूर्ण, चिन्हार्थमीमांसा हवी.
क्रिप्के यांनी दिलेली चिन्हार्थमीमांसा हे काम उत्तम रीतीने करते.

अभिजात तर्कशास्त्रासाठी केलेला खेळच येथेही खेळता येतो: चिन्हार्थमीमांसा
परिभाषित करायची आणि निर्दोषता व संपूर्णता सिद्ध करायची. मात्र पुढील भेद लक्षात
घेण्यासारखा आहे. अभिजात तर्कशास्त्रात चिन्हार्थमीमांसा एका अर्थाने
संयोजकांच्या अर्थातच अंतर्भूत असलेली “उघड” चिन्हार्थमीमांसा होती. क्रिप्के
चिन्हार्थमीमांसेसाठी काही अंतःप्रज्ञ स्पष्टीकरण देता येत असले, तरी तिच्यात
तशी अपरिहार्यता जाणवत नाही. शिवाय, अभिजात संरचना ही स्वाभाविक गणितीय
संकल्पना आहे; म्हणून संरचना ही संकल्पना अभिजात प्रथम-क्रम
तर्कशास्त्राच्या अभ्यासाचे साधन मानता येते, किंवा अभिजात प्रथम-क्रम
तर्कशास्त्र हे गणितीय संरचनेच्या अभ्यासाचे साधन मानता येते.
याउलट, क्रिप्के संरचना कडे फक्त तार्किक रचना म्हणूनच पाहता येते;
त्यांना स्वतंत्र गणितीय महत्त्व आहे असे दिसत नाही.

विधानीय भाषेसाठीची क्रिप्के संरचना~$\mModel M = \tuple{W, R, V}$
मध्ये $W$ हा संच, $W$ वरील लघुतम घटक असलेला $R$ हा अंशतः क्रम आणि
विधानीय चलांचे $W$ च्या घटक ना केलेले “एकस्वनिक” मूल्यांकन असते.
येथे $W$ चे घटक “जगे” किंवा “ज्ञानावस्था” दर्शवतात; $v \geq u$ हा
घटक $u$ ची “संभाव्य भावी अवस्था” दर्शवतो; आणि $u$ ला दिलेली विधानीय चरे ही
$u$ अवस्थेत सत्य असल्याचे माहीत असलेली विधाने असतात. बलन-संबंध
$\mSat{M}{\varphi}[w]$ मग हा संबंध भाषेतील कोणत्याही सूत्रे पर्यंत
विस्तारित करतो; $\mSat{M}{\varphi}[w]$ चा अर्थ “$w$ अवस्थेत $\varphi$ सत्य आहे” असा
वाचा. हा संबंध पुढीलप्रमाणे विगमनाने परिभाषित केला जातो:
\begin{enumerate}
\item $\mSat{M}{\Obj p_i}[w]$ तेव्हा आणि केवळ तेव्हाच, जेव्हा $\Obj p_i$
  हे $w$ ला दिलेल्या विधानीय चरांपैकी एक असते.
\item $\mSat/{M}{\lfalse}[w]$.
\item $\mSat{M}{(\varphi \land \psi)}[w]$ तेव्हा आणि केवळ तेव्हाच, जेव्हा
  $\mSat{M}{\varphi}[w]$ आणि $\mSat{M}{\psi}[w]$.
\item $\mSat{M}{(\varphi \lor \psi)}[w]$ तेव्हा आणि केवळ तेव्हाच, जेव्हा
  $\mSat{M}{\varphi}[w]$ किंवा $\mSat{M}{\psi}[w]$.
\item $\mSat{M}{(\varphi \lif \psi)}[w]$ तेव्हा आणि केवळ तेव्हाच, जेव्हा
  $w' \geq w$ आणि $\mSat{M}{\varphi}[w']$ असताना नेहमी $\mSat{M}{\psi}[w']$.
\end{enumerate}
प्रतिउदाहरण देणारी क्रिप्के संरचना रचून $\lnot (p \land q) \lif
(\lnot p \lor \lnot q)$ हे अंतःप्रज्ञावादी रीतीने वैध नाही हे दाखवण्याचा
प्रयत्न करणे हा चांगला सराव आहे.












\section{मोडल तर्कशास्त्रे}\label{fol:byd:mod:sec}

पुढील सशर्त वाक्याचे उदाहरण पाहा:
\begin{quote}
  जेरेमी त्या खोलीत एकटा असेल, तर त्याने मद्यपान केलेले आहे, तो विवस्त्र आहे
  आणि खुर्च्यांवर नाचत आहे.
\end{quote}
हे सशर्त विधान वास्तविक अभिव्यंजनाच्या अर्थाने सत्य असू शकते, पण तरीही
दिशाभूल करणारे असू शकते; कारण पूर्वांग असत्य आहे किंवा उत्तरांग सत्य आहे
एवढ्यापेक्षा पूर्वांग
आणि निष्कर्ष यांच्यात अधिक दृढ संबंध आहे असे ते सुचवताना दिसते. म्हणजे, हा
दावा फक्त या विशिष्ट जगातच सत्य नाही (जेरेमी खोलीत एकटा नसल्यामुळे येथे तो
क्षुल्लकपणे सत्य असू शकतो), तर पूर्वांग सत्य \emph{असते}, तर निष्कर्षही सत्य
\emph{असता}, असे त्याची शब्दरचना सुचवते. दुसऱ्या शब्दांत, हे विधान फक्त या
जगात नव्हे, तर कोणत्याही “संभाव्य” जगात सत्य आहे; किंवा ते या विशिष्ट जगात
केवळ सत्य असण्याऐवजी \emph{अनिवार्यपणे} सत्य आहे, असा त्याचा अर्थ घेता येतो.

अशा प्रकारच्या अनिवार्यतेचा अर्थ लावण्यासाठी मोडल तर्कशास्त्र रचले गेले.
नेहमीच्या विधानीय तर्कशास्त्रात चौकट-कारक जोडल्यावर मोडल विधानीय
तर्कशास्त्र मिळते; म्हणजे $\varphi$ हे सूत्र असेल, तर $\Box \varphi$ हेही
सूत्र असते. सहज अर्थाने, $\Box \varphi$ हे $\varphi$ \emph{अनिवार्यपणे} सत्य
आहे, किंवा कोणत्याही संभाव्य जगात सत्य आहे, असे सांगते. $\Diamond \varphi$ हे
सहसा $\lnot \Box \lnot \varphi$ चे संक्षिप्त रूप मानले जाते आणि $\varphi$
\emph{संभाव्यतः} सत्य आहे असे सांगणारे म्हणून वाचता येते. अर्थात, विधेय
तर्कशास्त्रातही मोडलता जोडता येते.

मोडल तर्कशास्त्राची चिन्हार्थमीमांसा देण्यासाठी क्रिप्के संरचना वापरता
येतात; खरे तर क्रिप्के यांनी ही चिन्हार्थमीमांसा प्रथम मोडल तर्कशास्त्र
डोळ्यांसमोर ठेवूनच रचली. अंशतः क्रमांपुरते मर्यादित न राहता, अधिक सामान्यपणे
“संभाव्य जगांचा” $P$ हा संच आणि जगांमधील $\Atom{R}{x,y}$ हा द्विपदी
“प्राप्यता” संबंध घेतला जातो. सहज अर्थाने, $\Atom{R}{p,q}$ हे जग~$q$ हे
जग~$p$ शी सुसंगत आहे असे सांगते; म्हणजे आपण जग~$p$ “मध्ये” असू, तर जग
$q$ सारखे असू शकले असते ही शक्यता विचारात घ्यावी लागते.

मोडल तर्कशास्त्राला कधी कधी “विस्तारात्मक” तर्कशास्त्राच्या विरुद्ध
“आशयात्मक” तर्कशास्त्र म्हणतात. अभिजात तर्कशास्त्रासारख्या विस्तारात्मक
तर्कशास्त्राची अभिप्रेत चिन्हार्थमीमांसा केवळ “प्रत्यक्ष” अशा एका जगाचा
निर्देश करते; पण “आशयात्मक” तर्कशास्त्राची चिन्हार्थमीमांसा अधिक विस्तृत
सत्ताशास्त्रावर अवलंबून असते. अनिवार्यतेची संरचना रचण्याखेरीज,
उपयोजनानुसार $\Box$ आणि $\Diamond$ यांचा नवा अर्थ लावून इतर भाषिक रचनांची
संरचना रचण्यासाठीही मोडलता वापरता येते. उदाहरणार्थ:
\begin{enumerate}
\item सिद्धतायोग्यता तर्कशास्त्रात $\Box \varphi$ चा अर्थ “$\varphi$ सिद्धतायोग्य
  आहे” आणि $\Diamond \varphi$ चा अर्थ “$\varphi$ सुसंगत आहे” असा घेतला जातो.
\item ज्ञानविषयक तर्कशास्त्रात $\Box \varphi$ चा अर्थ “मला $\varphi$ माहीत आहे” किंवा
  “माझा $\varphi$ वर विश्वास आहे” असा घेता येतो.
\item कालिक तर्कशास्त्रात $\Box \varphi$ चा अर्थ “$\varphi$ नेहमी सत्य आहे” आणि
  $\Diamond \varphi$ चा अर्थ “$\varphi$ कधीतरी सत्य असते” असा घेता येतो.
\end{enumerate}

मोडलता हाताळणारे नियम आणि स्वयंसिद्धके तर्कशास्त्रात जोडावीशी वाटतात.
उदाहरणार्थ, $\Log{S4}$ प्रणालीत विधानीय तर्कशास्त्राची नेहमीची स्वयंसिद्धके
आणि नियम, तसेच पुढील स्वयंसिद्धके असतात:
\begin{align*}
& \Box (\varphi \lif \psi) \lif (\Box \varphi \lif \Box \psi)\\
& \Box \varphi \lif \varphi \\
& \Box \varphi \lif \Box \Box \varphi
\intertext{आणि “$\varphi$ पासून $\Box \varphi$ हा निष्कर्ष काढा” हा नियमही असतो.
  $\Log{S5}$ मध्ये पुढील स्वयंसिद्धक जोडले जाते:}
& \Diamond \varphi \lif \Box \Diamond \varphi
\end{align*}
या स्वयंसिद्धकांच्या विविध आवृत्त्या वेगवेगळ्या उपयोजनांसाठी योग्य असू शकतात;
उदाहरणार्थ, S5 हे सहसा तार्किक अनिवार्यतेची संकल्पना लक्षणित करते असे मानतात.
चांगली गोष्ट अशी की क्रिप्के संरचना मधील प्राप्यता संबंध स्वाभाविक
रीतीने निर्बंधित करून, संबंधित निष्पत्ती प्रणाली निर्दोष आणि संपूर्ण
असेल अशी चिन्हार्थमीमांसा सहसा सापडते. उदाहरणार्थ, $\Log{S4}$ अशा क्रिप्के
संरचनेच्या वर्गाशी संबंधित आहे ज्यांत प्राप्यता संबंध परावर्ती आणि
संक्रमक असतो. $\Log{S5}$ अशा क्रिप्के संरचनेच्या वर्गाशी संबंधित
आहे ज्यांत प्राप्यता संबंध {\em वैश्विक} असतो, म्हणजे प्रत्येक जग इतर प्रत्येक
जगापासून प्राप्य असते; म्हणून $\Box \varphi$ तेव्हा आणि केवळ तेव्हाच लागू होते,
जेव्हा $\varphi$ प्रत्येक जगात लागू होते.












\section{इतर तर्कशास्त्रे}\label{fol:byd:oth:sec}

आतापर्यंत तुमच्या लक्षात आले असेलच की नवे तर्कशास्त्र रचणे अवघड नाही. तुम्हीही
स्वतःचा विन्यास तयार करू शकता, निगमन-पद्धती बनवू शकता आणि तिला साजेशी
चिन्हार्थमीमांसा घडवू शकता. निष्पत्ती प्रणाली चिन्हार्थमीमांसेसाठी
संपूर्ण हवी असेल, तर थोडे चातुर्य वापरावे लागू शकते; तसेच तुमचे तर्कशास्त्र
खरोखर रंजक आहे हे जगाला पटवून देण्यासाठी काही प्रयत्न करावे लागू शकतात. पण
त्याबदल्यात त्याच्या गणितीय आणि संगणकीय गुणधर्मांचा शोध घेताना अनेक तास
निर्भेळ आनंद मिळवता येईल.

अलीकडच्या दशकांत आकारिक तर्कशास्त्रांचा खरोखरच स्फोट झाला आहे. संदिग्ध
गुणधर्मांविषयीच्या तर्कविचाराचे प्रतिमान देण्यासाठी अस्पष्ट तर्कशास्त्र रचले
आहे. अनिश्चिततेविषयीच्या तर्कविचाराचे प्रतिमान देण्यासाठी संभाव्यताधारित
तर्कशास्त्र रचले आहे. खंडनीय प्रकारच्या तर्कविचाराचे प्रतिमान देण्यासाठी
पूर्वमान्य तर्कशास्त्रे आणि अ-एकस्वनिक तर्कशास्त्रे रचली आहेत; म्हणजे नवी
माहिती मिळाल्यावर नंतर उलटवता येतील असे “वाजवी” निष्कर्ष. ज्ञानाविषयीच्या
तर्कविचाराचे प्रतिमान देणारी ज्ञानविषयक तर्कशास्त्रे आहेत; कार्यकारण
संबंधांविषयीच्या तर्कविचाराचे प्रतिमान देणारी कार्यकारणविषयक तर्कशास्त्रे
आहेत; आणि नैतिक व नीतिविषयक कर्तव्यांविषयीच्या तर्कविचाराचे प्रतिमान देणारी
“कर्तव्यविषयक” तर्कशास्त्रेही आहेत. या प्रणाली मांडण्यामागील मुख्य प्रेरणा
तात्त्विक, गणितीय किंवा संगणकीय आहे यानुसार, अशा गोष्टींचा अभ्यास गणितीय
तर्कशास्त्र, तात्त्विक तर्कशास्त्र, कृत्रिम बुद्धिमत्ता, संज्ञानविज्ञान किंवा
इतर क्षेत्रांच्या नावाखाली झालेला आढळेल.

ही यादी सतत वाढत जाते आणि शक्यता अंतहीन दिसतात. सर्व मानवी तर्कबुद्धीचे गणनेत
रूपांतर करण्याचे लायप्निट्स यांचे स्वप्न आपण कदाचित कधीच साध्य करणार नाही---
पण त्यामुळे प्रयत्न करणे थांबवण्याचे कारण नाही.



